\chapter{Elaboration: Categorical and Analytic Foundations}
\label{chap:elaboration}

\section{Introduction and categorical context}
\label{sec:elaboration-intro}

This chapter collects the categorical, algebraic, and analytic details
that underpin the definitions of \hyperref[chap:qou]{Chapter}. The reader
who has read \hyperref[chap:qou]{Chapter} may return to specific sections as
needed; the remainder of the paper does not require this chapter in
linear order.

\begin{remark}[TQFT context]
  \label{rem:tqft-context}
  \uses{def:quantum-universe, def:quantum-observable-universe, def:c-quantum-observable-universe}
The axioms of
\hyperref[def:quantum-universe]{Definitions}
and [](#def:quantum-observable-universe) are grounded in the algebraic
framework of topological quantum field theory (TQFT).

\emph{Frobenius--Hopf objects and 2d TQFT.}
A commutative Frobenius algebra over $R$ is equivalent to a
$2$-dimensional TQFT
$Z \colon \mathbf{2Cob} \to R\text{-}\mathrm{Mod}$
\cite{kock2004, atiyah1988}. The Frobenius--Hopf object $A$ in
\hyperref[def:quantum-universe]{Definition} carries both the Frobenius structure
(encoding the $2$d cobordism data: the product $\mu$ corresponds to a
pair-of-pants cobordism, the Frobenius form $\lambda$ to a cap) and a
Hopf structure (encoding the quantum-group symmetry that extends the theory
to three dimensions). Concretely, a $2$d TQFT assigns to
the circle $S^1$ a Frobenius algebra $A$ and to each cobordism
$\Sigma \colon S^1 \sqcup \cdots \to S^1 \sqcup \cdots$ a linear map
built from $\mu$, $\Delta_\mu$, $\eta$, and $\lambda$. The Hopf
structure $(\Delta, \varepsilon, S)$ equips $A$ with the additional data
needed for three-dimensional and extended TQFTs
\cite{turaev2010, witten1988}.

\emph{The observability condition as a categorical contact condition.}
In the TQFT setting, the condition
$\mathrm{Hom}_{\mathcal{C}}(\mathbf{1}, \mathrm{Ad}(A)) = 0$
(\hyperref[def:quantum-observable-universe]{Definition}) is the categorical
analogue of the \emph{contact condition} (maximal non-integrability).
A non-zero morphism $\mathbf{1} \to \mathrm{Ad}(A)$ would give a
global invariant section of the adjoint object --- the categorical
counterpart of a nowhere-vanishing section of the adjoint bundle --- and
would force the TQFT to decompose into a direct sum of
lower-dimensional theories. Requiring
$\mathrm{Hom}(\mathbf{1}, \mathrm{Ad}(A)) = 0$ therefore ensures:

\begin{enumerate}
  \item \emph{irreducibility}: the associated $3$-dimensional TQFT
\end{enumerate}
   (Chern--Simons type, with gauge datum $A$ and quantum
   parameter $q$) does not split as a direct sum;
\begin{enumerate}
  \item \emph{non-degeneracy across cobordisms}: the TQFT pairing on
\end{enumerate}
   state spaces assigned to surfaces is non-degenerate
   \cite{freed2010};
\begin{enumerate}
  \item \emph{full dualisability}: in the sense of the cobordism
\end{enumerate}
   hypothesis \cite{baez1995, lurie2009}, the object $A$ is
   fully dualisable in the relevant Morita higher category, so the
   TQFT extends to a fully local (extended) field theory.

In the archimedean realisation
(\hyperref[def:c-quantum-observable-universe]{Example}), this categorical
condition reduces to
$H^0(M, \mathrm{Ad}(P)) = 0$ --- the classical requirement that the
instanton connection admits no infinitesimal automorphisms, which is the
gauge-theoretic manifestation of the contact (maximally non-integrable)
geometry of $(M, \theta)$.
\end{remark}

\begin{remark}[Vanishing conditions]
  \label{rem:vanishing-explication}
  \uses{def:c-quantum-observable-universe}
In \hyperref[def:c-quantum-observable-universe]{Example}, the condition
$H^0(M, \mathrm{Ad}(P)) = 0$ eliminates all infinitesimal gauge
symmetries, while $\Gamma_c(M,E) = 0$ ensures the jet-bundle
map $\Gamma(J^1 E) \hookrightarrow \mathcal{A}(P)$ is injective.
\end{remark}

\begin{remark}[Generalisation beyond strict pseudoconvexity]
  \label{rem:levi-generalisation}
  \uses{def:c-quantum-universe, def:c-quantum-observable-universe, def:q-harmonic-form, rem:c-substrate-via-levi}
\hyperref[def:c-quantum-universe]{Example} requires only Levi non-degeneracy
(not strict pseudoconvexity). This generality is deliberate: the QOU
framework applies whenever the following three cohomological/obstruction
conditions hold.

\begin{enumerate}
  \item \emph{Kohn--Rossi finite-dimensionality.}
\end{enumerate}
   The tangential CR complex
   $(\Omega^{0,\bullet}_b(M),\, \bar\partial_b)$ has
   finite-dimensional cohomology
   $H^{0,k}_{\bar\partial_b}(M)$ for all $k$. This replaces
   strict pseudoconvexity in ensuring analytic control over
   $q$-harmonic forms (\hyperref[def:q-harmonic-form]{Definition}).
   In the strictly pseudoconvex case, Kohn's sub-elliptic
   estimates guarantee this automatically; for indefinite
   Levi signature, one imposes it as a hypothesis.
\begin{enumerate}
  \item \emph{Non-degeneracy of the substrate.}
\end{enumerate}
   The Levi form $L_\theta$ is non-degenerate (but not necessarily
   positive definite), so the substrate parameter
   $q(p) = 1 + \lvert\langle L_\theta, F_A^+\rangle\rvert$
   (\hyperref[rem:c-substrate-via-levi]{Remark}) is well-defined and
   the regime $q > 1$ is attainable. When the Levi form has
   indefinite signature $(r,s)$ with $r,s \geq 1$, the Webster
   metric $g_\theta$ is pseudo-Riemannian; the contraction
   $\langle L_\theta, F_A^+\rangle$ remains well-defined (using
   the pseudo-Riemannian inner product on $\Lambda^+ H^*$), and
   one takes the absolute value to ensure $q \geq 1$.
\begin{enumerate}
  \item \emph{QOU vanishing (\hyperref[def:c-quantum-observable-universe]{Example}).}
\end{enumerate}
   The conditions $H^0(M, \mathrm{Ad}(P)) = 0$ and
   $\Gamma_c(M,E) = 0$ serve as the gauge-theoretic obstructions.
   These are purely topological/cohomological and independent of the
   Levi signature.

Together, conditions (1)--(3) ensure that all results of this chapter
(irreducibility of the connection, self-duality of the curvature,
substrate dynamics, $q$-deformed operators) extend to CR manifolds of
arbitrary Levi signature. The strictly pseudoconvex case (signature
$(n,0)$) is the simplest and motivates all examples in this chapter;
the indefinite case arises naturally for real hypersurfaces in
indefinite Kähler manifolds and for boundaries of
Grauert tubes over pseudo-Riemannian manifolds.
\end{remark}

\begin{remark}[Quantum gauge group: category-theoretic perspective]
  \label{rem:quantum-gauge-group-category}
  \uses{def:quantum-gauge-group}
The quantum gauge group
(\hyperref[def:quantum-gauge-group]{Definition}) is realised pointwise by
the Hopf algebra $\mathcal{O}_{q(p)}(\mathrm{SU}(2))$ of regular
functions on $\mathrm{SU}_{q(p)}(2)$. The categorical view replaces
each fibre with its representation category
$\mathrm{Rep}(\mathrm{SU}_{q(p)}(2))$, a braided monoidal category.
Tannaka--Krein duality reconstructs $\mathcal{O}_{q(p)}$ from
$\mathrm{Rep}(\mathrm{SU}_{q(p)}(2))$ via
$\mathrm{End}^{\otimes}(\omega)$, where $\omega$ is the forgetful
fibre functor. For any two points $p, p' \in M$ connected by a
path $\gamma$, the passage between fibres is captured by the
commutative square

\begin{equation}\label{eq:tannaka-krein-functoriality}
\begin{tikzcd}[row sep=2.5em, column sep=3.5em]
  \mathcal{O}_{q(p)}(\mathrm{SU}(2))
    \ar[r, "\mathrm{Rep}"]
    \ar[d, "\text{transition along } \gamma"']
  & \mathrm{Rep}(\mathrm{SU}_{q(p)}(2))
    \ar[d, "\text{monoidal functor}"]
  \\
  \mathcal{O}_{q(p')}(\mathrm{SU}(2))
    \ar[r, "\mathrm{Rep}"']
  & \mathrm{Rep}(\mathrm{SU}_{q(p')}(2))
\end{tikzcd}
\end{equation}

Tannaka--Krein reconstruction inverts the horizontal arrows, so the
Hopf-algebra and representation-category viewpoints are equivalent
fibrewise.
\end{remark}

\begin{remark}[Hopf-algebra structure: coassociativity]
  \label{rem:hopf-structure-diagrams}
  \uses{rem:hopf-non-cocommutativity, rem:self-dual-ses}
Write $\mathcal{O}_q = \mathcal{O}_q(\mathrm{SU}(2))$.
The Hopf comultiplication satisfies \emph{coassociativity}:

\begin{equation}\label{eq:coassociativity}
\begin{tikzcd}[row sep=2.5em, column sep=3.5em]
  \mathcal{O}_q
    \ar[r, "\Delta_q"]
    \ar[d, "\Delta_q"']
  & \mathcal{O}_q \otimes \mathcal{O}_q
    \ar[d, "\mathrm{id} \otimes \Delta_q"]
  \\
  \mathcal{O}_q \otimes \mathcal{O}_q
    \ar[r, "\Delta_q \otimes \mathrm{id}"']
  & \mathcal{O}_q \otimes \mathcal{O}_q \otimes \mathcal{O}_q
\end{tikzcd}
\end{equation}

Non-cocommutativity and the self-duality short exact sequence are
developed categorically in
\hyperref[rem:hopf-non-cocommutativity]{Remarks}
and [](#rem:self-dual-ses).
\end{remark}

\begin{remark}[Self-duality as constrained 2-group gauge theory]
  \label{rem:self-dual-2group}
  \uses{def:instanton}
A strict $2$-group is equivalently a \emph{crossed module}: a
group homomorphism $\partial \colon H \to G$ together with an
action of $G$ on $H$ satisfying the Peiffer identity. The
associated exact sequence

\begin{equation}\label{eq:crossed-module-ses}
\begin{tikzcd}[column sep=2em]
  1 \ar[r]
  & \ker\partial \ar[r]
  & H \ar[r, "\partial"]
  & G \ar[r]
  & \mathrm{coker}\,\partial \ar[r]
  & 1
\end{tikzcd}
\end{equation}

encodes the $2$-group as a homotopy type with
$\pi_0 = \mathrm{coker}\,\partial$ and
$\pi_1 = \ker\partial$.

The \emph{inner-automorphism $2$-group}
$\mathrm{INN}(G) = (G \xrightarrow{\mathrm{id}} G)$
has $\ker\partial = 1$ and $\mathrm{coker}\,\partial = 1$, so it is
connected and simply connected as a $2$-group. A $2$-connection for
$\mathrm{INN}(G)$ on a principal $2$-bundle over $M$ consists of a
$\mathfrak{g}$-valued $1$-form $A$ and a $\mathfrak{g}$-valued
$2$-form $B$, subject to the \emph{fake-flatness} condition

\[
F_A \;=\; B.$$

In the QOU setting, $G = \mathrm{SU}(2)$ and the Levi distribution
$H = \ker\theta$ has $\dim_{\mathbb{R}} H = 4$, so the Hodge
decomposition $\Lambda^2 H^* = \Lambda^+\!H^* \oplus \Lambda^-\!H^*$
is available. Imposing the *simplicity constraint*

B \;\in\; \Gamma\!\bigl(\Lambda^+\!H^* \otimes \mathrm{Ad}(P)\bigr)$$
\]

--- that is, requiring the $2$-form datum to be self-dual ---
and combining with fake flatness gives
the commutative diagram

\begin{equation}\label{eq:2group-instanton}
\begin{tikzcd}[row sep=2.5em, column sep=3em]
  \Omega^1(M,\,\mathfrak{g})
    \ar[r, "d_A"]
  & \Omega^2(M,\,\mathfrak{g})
    \ar[r, "\pi^+"]
  & \Omega^{2,+}(M,\,\mathfrak{g})
  \\
  & F_A \ar[u, "=\;B"]
    \ar[r, equal]
  & F_A^+ \ar[u, "\text{simplicity}"']
\end{tikzcd}
\end{equation}

The fake-flatness condition $F_A = B$ combined with the simplicity
constraint $B = F_A^+$ forces $F_A = F_A^+$, hence $F_A^- = 0$.
This is precisely the instanton equation
$F_A|_H = {*_H}\,F_A|_H$
of \hyperref[def:instanton]{Definition}. Thus, the self-dual formulation
used throughout this paper is a special case of $2$-group gauge theory
for $\mathrm{INN}(\mathrm{SU}(2))$ subject to a simplicity constraint
on the $2$-form datum.

This perspective clarifies the relationship between the instanton
condition in Chapter 1 and the higher-categorical structures of later
chapters: the parallel-transport $2$-functor of the
$\mathrm{INN}(\mathrm{SU}(2))$ $2$-bundle assigns to paths in $M$
not merely group elements in $\mathrm{SU}(2)$ but $2$-morphisms
encoding the self-dual curvature, with the simplicity constraint
selecting the physical (instanton) sector.
\end{remark}

\begin{proposition}[From Tannakian emergence to 2-group fake flatness]
  \label{prop:emergence-to-fake-flat}
  \uses{def:categorical-volume-constraint, thm:su2-emergence, rem:self-dual-2group, def:instanton}
Let $(\mathcal{C}, A, \omega)$ be a quantum observable universe satisfying
the categorical volume constraint
(\hyperref[def:categorical-volume-constraint]{Definition}), so that
$G_q = \mathrm{Aut}^{\otimes}(\omega) \cong \mathrm{SU}_q(2)$
(\hyperref[thm:su2-emergence]{Theorem}). Then the inner-automorphism
$2$-group $\mathrm{INN}(G_q)$ is the unique strict $2$-group whose
fake-flatness condition $F_A = B$ reproduces the
instanton equation $F_A^{-} = 0$ under the simplicity constraint.

---

The argument proceeds in three steps.

\emph{Step 1: The volume constraint determines the Lie algebra.}
The categorical volume constraint provides a non-degenerate epimorphism
$\epsilon \colon \Lambda^2 V \to \mathbf{1}$. Under the fibre functor
$\omega$, this endows the rank-$2$ module $\omega(V)$ with a
$q$-antisymmetric pairing whose invariance algebra is
$\mathfrak{su}_q(2)$. This is the Lie algebra of the reconstructed
group $G_q$, and hence of the crossed module
$\partial \colon G_q \xrightarrow{\mathrm{id}} G_q$ defining
$\mathrm{INN}(G_q)$.

*Step 2: The determinant-one condition descends to the
2-form datum.*
A $2$-connection $(A, B)$ for the principal $2$-bundle of
$\mathrm{INN}(G_q)$ takes values in
$\mathfrak{g}_q = \mathfrak{su}_q(2)$. Fake
flatness sets $B = F_A$, so $B$ inherits the
tracelessness forced by the determinant constraint
$\det(g_V) = 1$ of \hyperref[thm:su2-emergence]{Theorem}: $B$ is
$\mathrm{Ad}(P)$-valued with $\mathrm{tr}\,B = 0$, i.e.
$\mathfrak{su}_q(2)$-valued rather than $\mathfrak{gl}_q(2)$-valued.

*Step 3: Simplicity plus fake flatness gives the instanton
equation.*
Imposing the simplicity constraint
$B \in \Gamma(\Lambda^+\!H^* \otimes \mathrm{Ad}(P))$
(\hyperref[rem:self-dual-2group]{Remark}) and combining with $F_A = B$
yields $F_A = F_A^+$, hence $F_A^- = 0$. Step 1 guarantees the
structure group is $\mathrm{SU}_q(2)$---not a quotient of
$\mathrm{GL}_q(2)$---so the instanton moduli space carries the
gauge symmetry dictated by the categorical axioms.
\end{proposition}
\begin{proof}
The argument proceeds in three steps.

\emph{Step 1: The volume constraint determines the Lie algebra.}
The categorical volume constraint provides a non-degenerate epimorphism
$\epsilon \colon \Lambda^2 V \to \mathbf{1}$. Under the fibre functor
$\omega$, this endows the rank-$2$ module $\omega(V)$ with a
$q$-antisymmetric pairing whose invariance algebra is
$\mathfrak{su}_q(2)$. This is the Lie algebra of the reconstructed
group $G_q$, and hence of the crossed module
$\partial \colon G_q \xrightarrow{\mathrm{id}} G_q$ defining
$\mathrm{INN}(G_q)$.

*Step 2: The determinant-one condition descends to the
2-form datum.*
A $2$-connection $(A, B)$ for the principal $2$-bundle of
$\mathrm{INN}(G_q)$ takes values in
$\mathfrak{g}_q = \mathfrak{su}_q(2)$. Fake
flatness sets $B = F_A$, so $B$ inherits the
tracelessness forced by the determinant constraint
$\det(g_V) = 1$ of \hyperref[thm:su2-emergence]{Theorem}: $B$ is
$\mathrm{Ad}(P)$-valued with $\mathrm{tr}\,B = 0$, i.e.
$\mathfrak{su}_q(2)$-valued rather than $\mathfrak{gl}_q(2)$-valued.

*Step 3: Simplicity plus fake flatness gives the instanton
equation.*
Imposing the simplicity constraint
$B \in \Gamma(\Lambda^+\!H^* \otimes \mathrm{Ad}(P))$
(\hyperref[rem:self-dual-2group]{Remark}) and combining with $F_A = B$
yields $F_A = F_A^+$, hence $F_A^- = 0$. Step 1 guarantees the
structure group is $\mathrm{SU}_q(2)$---not a quotient of
$\mathrm{GL}_q(2)$---so the instanton moduli space carries the
gauge symmetry dictated by the categorical axioms.
\end{proof}

\begin{remark}[Categorical origin of the simplicity constraint]
  \label{rem:simplicity-from-category}
  \uses{thm:su2-emergence, rem:self-dual-2group}
The $*$-involution on $\mathcal{C}$---required by the QOU
observability axiom---induces via the fibre functor an involution on
$\Omega^2(M, \mathfrak{g}_q)$ whose $+1$-eigenspace is exactly the
self-dual forms $\Omega^{2,+}(M, \mathfrak{g}_q)$. Thus the full
chain

\[
\text{volume constraint}
  \;\xRightarrow{\text{Tannaka}}\;
  G_q = \mathrm{SU}_q(2)
  \;\xRightarrow{\mathrm{INN}}\;
  F_A = B
  \;\xRightarrow{*\text{-structure}}\;
  B = F_A^+
  \;\Longrightarrow\;
  F_A^- = 0
\]

is forced by the axioms of the quantum observable universe, with no
additional geometric input beyond the CR structure of $M$.
\end{remark}

\begin{remark}[Frobenius structure of the C gauge algebra]
  \label{rem:frobenius-simplification}
  \uses{def:quantum-universe, def:quantum-observable-universe, def:c-quantum-universe, def:c-quantum-observable-universe, def:instanton, lem:irreducible-connection, rem:q-torsion-existence}
The Hopf algebra $\mathcal{O}_q(\mathrm{SU}(2))$ carries a natural
\emph{Frobenius structure}: the Haar state
$h_q \colon \mathcal{O}_q \to \mathbb{C}$ induces a non-degenerate
pairing $\langle a, b \rangle = h_q(ab)$, making $\mathcal{O}_q$ into
a Frobenius algebra with comultiplication $\Delta_\mu$ dual to the
multiplication. This Frobenius comultiplication is compatible with the
Hopf comultiplication $\Delta_q$ of the coassociativity diagram and
satisfies the \emph{Frobenius relation}:

\begin{equation}\label{eq:frobenius-relation-elab}
\begin{tikzcd}[row sep=2.5em, column sep=3.5em]
  \mathcal{O}_q \otimes \mathcal{O}_q
    \ar[r, "\mu"]
    \ar[d, "\mathrm{id} \otimes \Delta_\mu"']
  & \mathcal{O}_q
    \ar[d, "\Delta_\mu"]
  \\
  \mathcal{O}_q \otimes \mathcal{O}_q \otimes \mathcal{O}_q
    \ar[r, "\mu \otimes \mathrm{id}"']
  & \mathcal{O}_q \otimes \mathcal{O}_q
\end{tikzcd}
\end{equation}

where $\mu$ is the algebra multiplication and $\Delta_\mu$ the
Frobenius comultiplication (dual to $\mu$ under the Haar pairing).
The non-degeneracy of $h_q$ is expressed by requiring the composite

\begin{equation}\label{eq:frobenius-nondegen}
\begin{tikzcd}[column sep=3em]
  \mathcal{O}_q \otimes \mathcal{O}_q
    \ar[r, "\mu"]
  & \mathcal{O}_q
    \ar[r, "h_q"]
  & \mathbb{C}
\end{tikzcd}
\end{equation}

to be a perfect pairing, i.e. the induced map
$\mathcal{O}_q \to \mathcal{O}_q^*$ is an isomorphism.

This Frobenius--Hopf structure on $\mathcal{O}_q(\mathrm{SU}(2))$ is
precisely the datum required by the general
\hyperref[def:quantum-universe]{Definition}: the Haar state $h_q$ realises
the Frobenius form $\lambda \colon A \to \mathbf{1}$, and the
non-degeneracy is the $\mathbb{C}$
instance of condition (1) in
\hyperref[def:quantum-observable-universe]{Definition}.
In the archimedean realisation (\hyperref[def:c-quantum-universe]{Example}),
the $5$-tuple $(M, \theta, G, P, E)$ and the two analytic vanishing
conditions of \hyperref[def:c-quantum-observable-universe]{Example}
reduce to a triple:

\begin{enumerate}
  \item a CR manifold $(M, \theta)$ of hypersurface type with
\end{enumerate}
   non-degenerate Levi form;
\begin{enumerate}
  \item a non-degenerate Frobenius--Hopf object $A$ in the monoidal
\end{enumerate}
   category of vector bundles over $M$, playing the role of the
   triple $(G, P, E)$; and
\begin{enumerate}
  \item an instanton connection on the underlying $A$-bundle
\end{enumerate}
   (\hyperref[def:instanton]{Definition}).

The vanishing condition $H^0(M, \mathrm{Ad}(P)) = 0$ then follows from
the non-degeneracy of the Frobenius form $\lambda$ via the diagram

\begin{equation}\label{eq:frobenius-vanishing}
\begin{tikzcd}[row sep=2.5em, column sep=2.5em]
  \mathrm{Ad}(A) \otimes \mathrm{Ad}(A)
    \ar[r, "\mu_{\mathrm{Ad}}"]
    \ar[d, "\text{global sections}"']
  & \mathrm{Ad}(A)
    \ar[r, "\lambda_{\mathrm{Ad}}"]
  & \mathcal{O}_M
    \ar[d, "\textstyle\int_M"]
  \\
  H^0(M,\,\mathrm{Ad}(A)) \otimes H^0(M,\,\mathrm{Ad}(A))
    \ar[rr, "\text{pairing}"']
  & & \mathbb{C}
\end{tikzcd}
\end{equation}

The adjoint representation $\mathrm{Ad}(A)$ inherits a non-degenerate
Frobenius pairing from $\lambda$; any global parallel section would
pair non-trivially with itself, contradicting the irreducibility of
the instanton connection
(\hyperref[lem:irreducible-connection]{Lemma}). Hence the Frobenius
non-degeneracy, combined with irreducibility, forces
$H^0(M, \mathrm{Ad}(A)) = 0$. The section-vanishing condition
$\Gamma_c(M, E) = 0$ follows similarly from the non-existence of
$A$-equivariant sections of the associated $\mathbb{P}^1(\mathbb{C})$-bundle.

Note that the horizontal distribution $H = \ker\theta$ itself
\emph{fails} to carry a Frobenius structure (as discussed in
\hyperref[rem:q-torsion-existence]{Remark}); it is the gauge algebra, not the
distribution, that is Frobenius.
\end{remark}

\begin{remark}[J^1E as the space of local states]
  \label{rem:j1e-local-states}
  \uses{lem:irreducible-connection, def:substrate-parameter, def:quantum-gauge-group}
The jet bundle $J^1 E$ serves as the space of local states of the
$\mathbb{C}$ quantum observable universe. The classical gauge dynamics,
encoded by the \emph{instanton connection} $A$
(\hyperref[lem:irreducible-connection]{Lemma}), are weighted by the substrate
parameter $q$ (\hyperref[def:substrate-parameter]{Definition}) to define the
quantum field algebra $\mathrm{SU}_q(2)$
(\hyperref[def:quantum-gauge-group]{Definition}).
\end{remark}

\begin{proposition}[d_q^2 does not vanish]
  \label{prop:dq-squared}
  \uses{def:q-ext-deriv}
Let $d_q$ be the $q$-deformed exterior derivative
(\hyperref[def:q-ext-deriv]{Definition}). For constant $q > 1$, the operator
$d_q^2$ is given by

\[
d_q^2\,\alpha
  \;=\; q^{-1}(1 - q^{-1})\,d\theta \wedge \mathcal{L}_\xi\,\alpha.
\]

In particular, $d_q^2 = 0$ if and only if $q = 1$ (classical case) or
$\alpha$ lies in the kernel of $\mathcal{L}_\xi$ (Reeb-invariant forms).
On horizontal $\mathcal{L}_\xi$-invariant forms, $d_q$ \emph{does}
define a cochain complex.

---

Write $a = q^{-1}$ and $b = 1 - q^{-1}$, so $a + b = 1$. Then

\begin{align*}
  d_q^2\,\alpha
  &= d_q\bigl(a\,d\alpha + b\,\theta \wedge \mathcal{L}_\xi\alpha\bigr) \\
  &= a^2\,d^2\alpha
     + ab\,d(\theta \wedge \mathcal{L}_\xi\alpha)
     + ab\,\theta \wedge \mathcal{L}_\xi\,d\alpha
     + b^2\,\theta \wedge \mathcal{L}_\xi(\theta \wedge \mathcal{L}_\xi\alpha).
\end{align*}

The first term vanishes since $d^2 = 0$. The second expands as
$ab\,d\theta \wedge \mathcal{L}_\xi\alpha
\begin{itemize}
  \item ab\,\theta \wedge d\mathcal{L}_\xi\alpha$.
\end{itemize}
The third term cancels the second's negative part because
$[d,\, \mathcal{L}_\xi] = 0$ implies
$\theta \wedge \mathcal{L}_\xi\,d\alpha
= \theta \wedge d\,\mathcal{L}_\xi\alpha$.
For the fourth term, the Reeb field satisfies
$\mathcal{L}_\xi\theta = 0$ (since $\iota_\xi d\theta = 0$ and
$\theta(\xi)=1$), so
$\mathcal{L}_\xi(\theta \wedge \mathcal{L}_\xi\alpha)
= \theta \wedge \mathcal{L}_\xi^2\alpha$;
the factor $\theta \wedge \theta = 0$ in $b^2\,\theta \wedge \theta
\wedge \mathcal{L}_\xi^2\alpha$ kills this term. Hence
$d_q^2\alpha = ab\,d\theta \wedge \mathcal{L}_\xi\alpha$.
\end{proposition}
\begin{proof}
Write $a = q^{-1}$ and $b = 1 - q^{-1}$, so $a + b = 1$. Then

\begin{align*}
  d_q^2\,\alpha
  &= d_q\bigl(a\,d\alpha + b\,\theta \wedge \mathcal{L}_\xi\alpha\bigr) \\
  &= a^2\,d^2\alpha
     + ab\,d(\theta \wedge \mathcal{L}_\xi\alpha)
     + ab\,\theta \wedge \mathcal{L}_\xi\,d\alpha
     + b^2\,\theta \wedge \mathcal{L}_\xi(\theta \wedge \mathcal{L}_\xi\alpha).
\end{align*}

The first term vanishes since $d^2 = 0$. The second expands as
$ab\,d\theta \wedge \mathcal{L}_\xi\alpha
\begin{itemize}
  \item ab\,\theta \wedge d\mathcal{L}_\xi\alpha$.
\end{itemize}
The third term cancels the second's negative part because
$[d,\, \mathcal{L}_\xi] = 0$ implies
$\theta \wedge \mathcal{L}_\xi\,d\alpha
= \theta \wedge d\,\mathcal{L}_\xi\alpha$.
For the fourth term, the Reeb field satisfies
$\mathcal{L}_\xi\theta = 0$ (since $\iota_\xi d\theta = 0$ and
$\theta(\xi)=1$), so
$\mathcal{L}_\xi(\theta \wedge \mathcal{L}_\xi\alpha)
= \theta \wedge \mathcal{L}_\xi^2\alpha$;
the factor $\theta \wedge \theta = 0$ in $b^2\,\theta \wedge \theta
\wedge \mathcal{L}_\xi^2\alpha$ kills this term. Hence
$d_q^2\alpha = ab\,d\theta \wedge \mathcal{L}_\xi\alpha$.
\end{proof}

\begin{proposition}[Jacobiator of the q-Lie bracket]
  \label{prop:q-jacobiator}
  \uses{def:q-lie-bracket}
Let $U, V, W \in \Gamma(H)$ be horizontal vector fields (i.e.
$\theta(U) = \theta(V) = \theta(W) = 0$). The \emph{Jacobiator}

\[
\mathrm{Jac}_q(U,V,W)
  \;:=\; [[U,V]_q, W]_q + [[V,W]_q, U]_q + [[W,U]_q, V]_q
\]

is given by

\[
\mathrm{Jac}_q(U,V,W)
  \;=\; -q^{-1}(1-q^{-1})\!\sum_{\mathrm{cyc}}
        d\theta(V,W)\,\mathcal{L}_\xi U,
\]

where the sum is over cyclic permutations of $(U,V,W)$. In particular,
$\mathrm{Jac}_q = 0$ if and only if $q = 1$ (classical Jacobi identity)
or $d\theta|_{H \times H} = 0$ (integrable horizontal distribution).
Since $d\theta|_{H \times H}$ is the Levi form $L_\theta$, which is
non-degenerate by hypothesis, $[\,\cdot\,,\cdot\,]_q$ violates the Jacobi
identity for all $q \neq 1$. Thus $(\mathfrak{X}(M), [\,\cdot\,,\cdot\,]_q)$
is a \emph{Leibniz algebra} (or a pre-Lie algebra), not a Lie algebra.

---

For horizontal $V, W$, the bracket reduces to $[V,W]_q = q^{-1}[V,W]$.
The mixed term $\theta([V,W]_q) = -q^{-1}\,d\theta(V,W)$ is in general
non-zero. Substituting into the next bracket:

\begin{align*}
  [[V,W]_q,\, U]_q
  &= q^{-2}\,[[V,W],U]
    - q^{-1}(1-q^{-1})\,d\theta(V,W)\,\mathcal{L}_\xi U,
\end{align*}

using $\theta(U)=0$ and $\theta([V,W]_q) = -q^{-1}d\theta(V,W)$.
Summing cyclically, the $q^{-2}$ terms assemble into $q^{-2}\,\mathrm{Jac}(U,V,W) = 0$
by the classical Jacobi identity, leaving only the $d\theta$ terms.
\end{proposition}
\begin{proof}
For horizontal $V, W$, the bracket reduces to $[V,W]_q = q^{-1}[V,W]$.
The mixed term $\theta([V,W]_q) = -q^{-1}\,d\theta(V,W)$ is in general
non-zero. Substituting into the next bracket:

\begin{align*}
  [[V,W]_q,\, U]_q
  &= q^{-2}\,[[V,W],U]
    - q^{-1}(1-q^{-1})\,d\theta(V,W)\,\mathcal{L}_\xi U,
\end{align*}

using $\theta(U)=0$ and $\theta([V,W]_q) = -q^{-1}d\theta(V,W)$.
Summing cyclically, the $q^{-2}$ terms assemble into $q^{-2}\,\mathrm{Jac}(U,V,W) = 0$
by the classical Jacobi identity, leaving only the $d\theta$ terms.
\end{proof}

\begin{remark}[Classical recovery of q-deformed operators]
  \label{rem:classical-recovery}
  \uses{def:q-ext-deriv, def:q-lie-bracket, def:q-connection, def:q-reeb-vector-field}
All four $q$-deformed operators interpolate between the classical ($q = 1$)
and deformed ($q > 1$) regimes via the uniform pattern
$q^{-1} \cdot (\text{classical}) + (1 - q^{-1}) \cdot (\text{Reeb correction})$.
At $q = 1$ each recovers its classical counterpart:
$d_q = d$ (\hyperref[def:q-ext-deriv]{Definition}),
$[\,\cdot\,,\cdot\,]_q = [\,\cdot\,,\cdot\,]$
(\hyperref[def:q-lie-bracket]{Definition}),
$A_q = A$ (\hyperref[def:q-connection]{Definition}), and
$\xi_q = \xi$ (\hyperref[def:q-reeb-vector-field]{Definition}).
The connection $A_q$ interpolates between the full instanton connection and
its Reeb component $\iota_\xi A\;\theta$.
\end{remark}

\begin{remark}[q-Torsion leaves as energy minima]
  \label{rem:q-torsion-leaf-minimum}
  \uses{def:q-connection, def:q-lie-bracket, def:q-torsion-leaf}
Let $A_q$ denote the $q$-deformed connection
(\hyperref[def:q-connection]{Definition}) and define the $q$-curvature

\[
\Omega_q \;=\; dA_q + \tfrac{1}{2}\,[A_q,\, A_q]_q,
\]

where $[\,\cdot\,,\cdot\,]_q$ is the $q$-deformed Lie bracket
(\hyperref[def:q-lie-bracket]{Definition}) extended to $\mathrm{Ad}(P)$-valued
forms. A $q$-torsion leaf $L$ (\hyperref[def:q-torsion-leaf]{Definition}) is a
local minimum of $\lVert\Omega_q\rVert$ restricted to the horizontal
distribution: the invariance of $L$ under $\phi_t^q$ forces the first
variation of $\lVert\Omega_q|_L\rVert$ to vanish, so $L$ sits in a
$q$-curvature well. Physically, a $q$-torsion leaf represents a region of
the CR manifold where the gauge curvature is minimised and the $q$-deformed
dynamics are in equilibrium.
\end{remark}

\begin{remark}[Physical interpretation of q-harmonicity]
  \label{rem:q-harmonic-physical}
  \uses{rem:reeb-physical, def:q-harmonic-form}
See \hyperref[rem:reeb-physical]{Remark} for the physical interpretation of
$q$-harmonic forms and the Reeb decomposition. In brief: the
$q$-Laplacian $\Delta_q$ interpolates between the Hodge Laplacian
($q = 1$) and the deep-quantum regime ($q \to \infty$), so a
$q$-harmonic form --- a "particle" --- is itself $q$-dependent.
\end{remark}

\begin{remark}[q-Deformed Schrodinger equation]
  \label{rem:schrodinger-jet-coaction}
  \uses{def:q-reeb-vector-field, def:q-harmonic-form, def:quantum-planck-constant, rem:heisenberg-commutator, prop:main-commutative-diagram, def:quantum-universe, def:quantum-observable-universe}
The $q$-deformed Schrodinger equation governs the evolution of states
in the linearised state bundle. Throughout this paper the symbol $i$ in
the Schrodinger equation is identified with the complex structure
operator $J \colon H \to H$ on the Levi distribution ($J^2 = -\mathrm{Id}$):
it acts as a geometric twisting operator on horizontal planes rather than
as an algebraic imaginary unit.

\emph{Linearisation convention.}
The fibre $\mathbb{P}^1(\mathbb{C})$ is a non-linear space, so one cannot write a
linear PDE for sections of $E$ directly. We work instead with the
\emph{linearised} model: identify a neighbourhood of a base-point in
$\mathbb{P}^1(\mathbb{C})$ with $\mathbb{C}$ via stereographic projection, and
interpret $\psi$ as a section of the line bundle $E_{\mathrm{lin}}$
obtained by linearising $E$ at the chosen section. The covariant
derivative $D_A = d + A$ then acts on sections of this line bundle in the
standard way.

For a section $\psi \in \Gamma(E_{\mathrm{lin}})$ and the $q$-deformed Reeb
vector field $\xi_q$ (\hyperref[def:q-reeb-vector-field]{Definition}), the
Schrodinger dynamics take the form

\[
J\,\hbar_q(p)\,\iota_{\xi_q} D_A \psi \;=\; \hat{H}_q\, \psi,
\]

where $\hat{H}_q = -\tfrac{1}{2}\hbar_q\,\Delta_q + V_q$ is the
\emph{$q$-deformed Hamiltonian} (here $\Delta_q$ is the $q$-Laplacian of
\hyperref[def:q-harmonic-form]{Definition} acting on sections of
$E_{\mathrm{lin}}$, and $V_q$ is a potential depending on the curvature
$\Omega_q$), and:

\begin{itemize}
  \item $\hbar_q(p) = 1 - q(p)^{-1}$: the quantum Planck constant
\end{itemize}
  (\hyperref[def:quantum-planck-constant]{Definition}) is a local field;
\begin{itemize}
  \item $D_A = d + A$: the instanton connection on $P$ acts as the potential;
  \item $\iota_{\xi_q} D_A \psi$: the covariant derivative $D_A\psi$ is an
\end{itemize}
  $E_{\mathrm{lin}}$-valued $1$-form; contracting it with $\xi_q$
  yields a section of $E_{\mathrm{lin}}$, matching the right-hand side.

As $q(p) \to 1$ (classicality), the commutator
$[\hat{X}, \hat{P}]_q \to 0$ (\hyperref[rem:heisenberg-commutator]{Remark}), the
uncertainty vanishes, and the wave-like nature of the section freezes into a
classical trajectory.

\emph{Categorical role.}
Although this remark is stated as a PDE rather than in categorical language,
it is essential to the paper's structure: the Schrodinger equation
provides the bottom row of the main commutative diagram
(\hyperref[prop:main-commutative-diagram]{Proposition}). Categorically, the
$q$-deformed Hamiltonian $\hat{H}_q$ defines an endomorphism of the
module category $\mathrm{Mod}_{E_{\mathrm{lin}}}$ over the jet algebra
$\Gamma(J^1 E)$, and the Schrodinger equation is the condition that
$\psi$ lies in the kernel of $\hat{H}_q - J\hbar_q\,\iota_{\xi_q}D_A$.
In the base-field-independent setting of
\hyperref[def:quantum-universe]{Definitions}
and [](#def:quantum-observable-universe), this extends to arbitrary base
rings $R$ by replacing $\Gamma(E_{\mathrm{lin}})$ with the appropriate
$R$-linear category of sections.
\end{remark}

\section{Lifting: jet prolongation, transport, and torsion}
\label{sec:lifting}

The categorical jet object $J^1(A)$ (\hyperref[def:categorical-jet]{Definition})
encodes first-order deformations of the Frobenius--Hopf object. In this
section we develop the \emph{lifting} programme: promoting a quantum
observable universe on $A$ to one on $J^1(A)$, and iterating to build
the brane tower.

\begin{proposition}[Jet prolongation]
  \label{prop:categorical-jet-prolongation}
  \lean{QOU.categorical_jet_prolongation}
  \uses{def:categorical-connection, def:categorical-jet}
Let $(\mathcal{C}, A, \omega)$ be a quantum universe with
connection $\nabla$ (\hyperref[def:categorical-connection]{Definition}).
Then $J^1(A)$ (\hyperref[def:categorical-jet]{Definition}) inherits a
Frobenius--Hopf structure from $A$, and the quadruple
$(\mathcal{C},\, J^1(A),\, \omega',\, q)$ --- where $\omega'$ is
the induced fibre functor --- is again a quantum universe.
The following diagram commutes:

\begin{equation}\label{eq:jet-prolongation-cd}
\begin{tikzcd}[row sep=2.5em, column sep=3em]
  J^1(A) \ar[r, "\nabla'"] \ar[d, "j^1"']
  & J^1(A) \otimes \Omega^1(J^1(A)) \ar[d, "j^1 \otimes j^1"] \\
  A \ar[r, "\nabla"']
  & A \otimes \Omega^1(A)
\end{tikzcd}
\end{equation}

where $\nabla'$ is the prolonged connection and $\nabla$ the original.
\end{proposition}

% ERROR: block "categorical-transport" not found
% ERROR: block "categorical-q-reeb-flow" not found
% ERROR: block "categorical-q-torsion-leaf" not found